\documentclass[11pt,letterpaper]{article}
\usepackage[T1]{fontenc}
\usepackage{lmodern}
\usepackage[margin=25mm]{geometry}
\usepackage[hidelinks]{hyperref}
\usepackage{xurl}
\pagestyle{empty}
\setlength{\parindent}{0pt}
\setlength{\parskip}{.65em}
\hypersetup{pdftitle={Cover letter: Adaptive HACT},pdfauthor={Da Xu and Fei Wang}}
\begin{document}
\textbf{Da Xu}\\
China Mobile Research Institute\\
\href{mailto:xudayj@chinamobile.com}{xudayj@chinamobile.com}
\hfill 23 October 2026

Editors\\
\emph{IEEE Transactions on Software Engineering}

Dear Editors,

Please consider \emph{Adaptive HACT: Information Limits and Safe Adaptation
of Verification Evidence}, by Da Xu and Fei Wang, for review as a regular
research article.

The paper studies how software-verification evidence can be reorganized
without changing its acceptance contract. Its principal contribution is an
explicit family of completion laws with identical exact singleton and
pairwise summaries but disjoint globally optimal evidence trees, together
with sharp minimax regret. A complete-epoch construction and fixed-eight
candidate-class extension connect that information limit to batched check
execution while keeping their distinct scopes explicit.

A complementary result gives an exact global baseline under exchangeable
complete batches. We attribute the cardinality-state principle to established
batch-update tree optimization and explain the model-specific specialization.
The paper also gives a conditional local publication invariant and a
paired-cost adoption rule that accounts for calibration, declared workload
shift and installation.

The empirical contribution is a reproducible boundary study. Retrospective
scoring preserves the symmetric baseline's item/file policy reversal and
original-holdout losses, the maintained-counter ledger's service advantage,
Packaging's resource exclusion and six initial VOID runs. The article makes
no general service or agent-speedup claim. This combination of information
limits, acceptance-preserving composition and independently checked workload
boundaries is intended to inform the design of verification infrastructure
for software agents.

The accompanying technical supplement provides full complete-epoch and
padding proofs and an artifact guide. Code, exact references, independent
audits and all retained observations are publicly available in the
versioned artifact. AI assistance is disclosed in the manuscript and
supplement. The submission metadata records the public repository versions.

Thank you for considering the work.

Sincerely,\\
Da Xu\\
Corresponding author for Da Xu and Fei Wang
\end{document}
